% chB 记号总表 (照录原书 List of main symbols, 书页305-309 / p323-327)
\chapter{记号总表}\markright{记号总表}

本表照录原书卷末“主要符号表”（List of main symbols）并译出含义：
每章只列该章新引入的记号，全书记号体系前后一致。建议在阅读
第 7--10 章时对照使用。

\section*{第 2--3 章（金融学导论与衍生证券）}

\begin{center}
\begin{tabular}{@{}L{3.2cm}L{10.6cm}@{}}
\toprule
$R$ & 组合收益 \\
$\equiv$ & 定义；恒等 \\
$E[X]$ & 随机变量 $X$ 的期望值 \\
$\sigma$ & 股票或组合的波动率 \\
$E[X_{n+1}|x_1,\dots,x_n]$ & 以 $x_1,\dots,x_n$ 为条件的 $X_{n+1}$ 期望 \\
$D(S)$ & 标的 $S$ 的衍生品价值 \\
$\Pi(S,t)$ & 股票价格 $S$ 的投资组合 \\
$r(t)$ & 即期利率 \\
$P(t,T)$ & $t$ 时刻的、$T$ 到期的国债价值 \\
$\mathcal{B}(t,T)$ & 附息国债价值 \\
$f(t,x)$ & 远期利率（$t$ 时刻约定的、未来 $x$ 时刻瞬时贷款利率） \\
\midrule
$F(t,T)$ & 远期合约价值（$t$ 时刻、$T$ 到期） \\
$\mathcal{F}(t,T)$ & 期货合约价值 \\
$C(t,S(t))$ & 看涨期权在 $t$ 时刻的价值 \\
$P(t)$ & 看跌期权价值 \\
$S(t)$ & $t$ 时刻股票价格 \\
$R(t)$ & 白噪声 \\
$P(r_n)$ & 离散化白噪声的概率分布 \\
$\epsilon$ & 无穷小量 \\
$x(t)$ & 股票价格的对数 $\ln S(t)$ \\
$\tau=T-t$ & 距到期时间 \\
$P_m(x)$ & Black--Scholes 鞅概率分布 \\
$N(d_{\pm})$ & 正态随机变量的累积分布 \\
$V=\sigma^2$ & 股票价格方差 \\
$\rho$ & 白噪声之间的相关系数 \\
$\lambda,\mu,\xi$ & 随机波动率参数 \\
\bottomrule
\end{tabular}
\end{center}

\section*{第 4--5 章（哈密顿量与路径积分）}

\begin{center}
\begin{tabular}{@{}L{3.2cm}L{10.6cm}@{}}
\toprule
$\mathcal{V}$ & 线性矢量空间（态空间） \\
$\mathcal{V}_{\mathrm{dual}}$ & 对偶线性矢量空间 \\
$\vert g\rangle$ & 属于 $\mathcal{V}$ 的右矢 \\
$\langle p\vert$ & 属于 $\mathcal{V}_{\mathrm{dual}}$ 的左矢 \\
$\vert x\rangle$ & 实直线上的右矢基矢 \\
$\psi(x)$ & 右矢 $\vert\psi\rangle$ 的分量 $\langle x\vert\psi\rangle$ \\
$H$ & 哈密顿量 \\
$\vert\psi_E\rangle$ & $H$ 的本征函数 \\
$E$ & $\vert\psi_E\rangle$ 的本征值 \\
$H_{BS}$, $H_{MG}$ & Black--Scholes / Merton--Garman 哈密顿量 \\
$p(x,y,T-t,x',y')$ & 随机波动率股票的定价核 \\
$\vert C,t\rangle$ & 期权价格右矢 \\
$C(t,x)$ & $\langle x\vert C,t\rangle$ \\
$p_{BS}(x,T-t,x')$ & Black--Scholes 定价核 \\
$\vert p\rangle$ & 动量右矢基矢 \\
$H_{\mathrm{eff}}$ & 有效期权定价哈密顿量 \\
$V(x)$ & 表示路径依赖期权的势 \\
$H_V$ & 含势哈密顿量 \\
$H_{DO}$, $H_{DB}$ & 障碍期权哈密顿量 \\
$\langle x\vert\phi_n\rangle$ & 双障碍期权的本征函数 \\
\midrule
$L$ & 拉格朗日量 \\
$S$ & 作用量 \\
$\int DX$, $\int DY$ & 路径积分测度 \\
$L_{BS}$ & Black--Scholes 拉格朗日量 \\
$g[x,K]$ & 路径依赖收益函数 \\
$Z(t_i,t_f;j)$ & 简谐振子的生成泛函 \\
$D(t,t';t_i,t_f)$ & 传播子 \\
$L_{MG}$ & Merton--Garman 拉格朗日量 \\
$S_0$, $S_1$ & 随机波动率作用量 \\
$Z(i,y,p)$ & 股票与波动率矩的生成函数 \\
$F$, $F'$ & 生成函数 $Z(i,y,p)$ 的解 \\
\bottomrule
\end{tabular}
\end{center}

\section*{第 6 章（随机利率）}

\begin{center}
\begin{tabular}{@{}L{3.2cm}L{10.6cm}@{}}
\toprule
$a(r,t)$, $\sigma(r,t)$ & 即期利率的漂移与波动率 \\
$P_F(r,t:r_0)$ & $r$ 的前向条件概率 \\
$H_F$, $H_B$ & 前向与倒向哈密顿量 \\
$L_F$, $S_F$ & 前向拉格朗日量与作用量 \\
$Z_B$, $S_B$ & 倒向配分函数与作用量 \\
$S_V$ & Vasicek 作用量 \\
$x$ & 未来时间 \\
$\alpha(t,x)$ & 远期利率的漂移速度 \\
$\sigma_i(t,x)$ & 远期利率的波动率函数 \\
$\mathbf{Z}[j,t_1,t_2]$ & HJM 模型的生成泛函 \\
$t_0$ & 当前时刻 \\
$t_*$ & 未来时刻（$>t_0$，通常为期权到期时刻） \\
$E_{[t_0,t_*]}[X]$ & $[t_0,t_*]$ 上全体随机变量的期望 \\
$T$, $\Delta_0$, $\mathcal{R}$ & 远期利率的积分区域 \\
$\mathcal{F}(t_0,t_*,T)$ & 国债期货价格 \\
$\Omega$, $\Omega_{\mathcal{F}}$ & 国债期货价格中的函数 \\
$C(t_0,t_*,T,K)$ & 国债看涨期权（到期 $t_*$、标的到期 $T$） \\
$q^2$ & 国债期权中的波动率 \\
$\Psi(G,t_*,T)$ & 国债期权价格中出现的函数 \\
$S_F^{bk}$, $S_B^{bk}$ & Black--Karasinski 前向与倒向作用量 \\
$L_B^{bk}$, $H_B^{bk}$ & Black--Karasinski 倒向拉格朗日量与哈密顿量 \\
\bottomrule
\end{tabular}
\end{center}

\section*{第 7 章（远期利率场论）}

\begin{center}
\begin{tabular}{@{}L{3.2cm}L{10.6cm}@{}}
\toprule
$S_{\mathrm{String}}$ & 弦作用量 \\
$Z$ & 远期利率配分函数 \\
$T_{FR}$ & 未来时间最大值 \\
$S[f]$, $\mathcal{L}[f]$ & 远期利率作用量与拉格朗日量 \\
$\int\mathcal{D}f$ & 对全部 $f(t,x)$ 的泛函积分 \\
$A(t,x)$ & 远期利率速度量子场 \\
$\mathbf{Z}[J]$ & 速度场的生成泛函 \\
$\delta/\delta J(t,x)$ & 泛函导数 \\
$\mu$ & 远期利率拉格朗日量中的刚性参数 \\
$D(x,x';t)$ & 常数刚性远期利率传播子 \\
$\theta$ & 未来时间坐标 $x-t$ \\
$\theta_{\pm}$ & 未来时间坐标 $\theta\pm\theta'$ \\
$m$ & 传播子对角轴的斜率 \\
$B(t_0,t_*)$ & 货币市场账户计价物 \\
$S_*$ & 计价物 $P(t,T_*)$ 下的远期利率作用量 \\
$\alpha_*(t,x)$ & 计价物 $P(t,T_*)$ 下的远期利率漂移 \\
$\phi(t,x)$ & （非线性）远期利率的对数 $\ln(f(t,x)/f_0)$ \\
$\psi[t,x;A]$ & 以 $A(t,x)$ 场表出的 $\phi(t,x)$ \\
$\mathcal{J}(t,x;t'x';A)$ & $\phi(t,x)\to A(t,x)$ 变量替换的矩阵 \\
$c^{\dagger}(t,x)$, $c(t,x)$ & 费米子量子场 \\
$\nu$ & 远期利率随机波动率指数 \\
$\rho$, $\xi$, $\kappa$ & 远期利率随机波动率参数 \\
$D_C(\theta,\theta')$ & 约束远期利率传播子 \\
$D_M(\theta,\theta')$ & 非常数刚性远期利率传播子 \\
$\lambda$ & 远期利率作用量中的刚度参数 \\
$G(\theta,\theta')$ & 远期利率刚度传播子 \\
$r_Q$ & 垂直于传播子对角轴的曲率 \\
$z(\theta)$ & 心理未来时间 \\
$Z_f[J]$ & 远期利率生成泛函 \\
$f_{mn}$, $A_{mn}$ & 格点远期利率与速度量子场 \\
\bottomrule
\end{tabular}
\end{center}

\section*{第 8 章（实证远期利率）}

\begin{center}
\begin{tabular}{@{}L{3.2cm}L{10.6cm}@{}}
\toprule
$L(t,T)$ & Libor 单利率 \\
$\tau$ & 90 天时间间隔 \\
$\langle\delta f(t,\theta)\rangle$ & 远期利率平均变化率 \\
$\langle\delta f(t,\theta)\delta f(t,\theta')\rangle_c$ & 远期利率速度关联函数 \\
$\kappa(t,\theta)$ & 远期利率峰度 \\
$C(\theta,\theta')$ & 归一化关联函数 \\
$\sigma_r$ & 即期利率波动率 \\
$C_R(\theta,\theta')$ & 常数刚性关联函数 \\
$C_C(\theta,\theta')$ & 约束即期利率关联函数 \\
$C_M(\theta,\theta')$ & 非常数刚性关联函数 \\
$C_Q(\theta,\theta')$ & 刚度传播子关联函数 \\
$C_{Qz}(\theta,\theta')$ & 刚度传播子+心理时间关联函数 \\
$R_Q$ & 垂直于归一化传播子对角轴的曲率 \\
$z_{\pm}$ & 心理未来时间的组合 $z(\theta)\pm z(\theta')$ \\
$\eta$ & 心理未来时间的标度指数 \\
\bottomrule
\end{tabular}
\end{center}

\section*{第 9 章（国债衍生品与对冲）}

\begin{center}
\begin{tabular}{@{}L{3.2cm}L{10.6cm}@{}}
\toprule
$\Delta$, $\Gamma$, $\Theta$ & 国债期权对冲参数（希腊字母） \\
$\mathcal{V}$, $\mathcal{V}_i$, $\mathcal{V}_{\sigma}$, $\mathcal{V}_{\mu}$ & 国债期权方差的对冲参数 \\
$\mathrm{Cap}(t_0,t_*,\tau,X)$ & 利率上限期权 \\
$L_*(t_0,t_*,\tau)$ & Libor $L(t_*,t_*,\tau)$ \\
$q_{\mathrm{cap}}$ & 利率上限期权波动率 \\
$\mathrm{Cap}_B(t_0,t_*,\tau,X)$ & 利率上限期权的 Black 公式 \\
$F_{\mathrm{cap}}$ & 利率上限期权的远期合约 \\
$n_1(t)$ & 债券 Delta 对冲的对冲权重 \\
$B(r;t;T)$ & 给定 $r$ 时国债的条件期望 \\
$\Delta_i$ & 方差最小化对冲权重 \\
$\mathrm{Var}[\Pi(t_*)]$ & 未来时刻 $t_*$ 债券组合方差 \\
$V_0$ & （未对冲）国债方差 \\
$V$ & 对冲组合残余方差最小值 \\
$V_*$ & 有限时间对冲组合残余方差最小值 \\
$P_i$ & 等于 $P(t,T_i)$ \\
$M_{ij}$ & 国债对冲矩阵 \\
$L_i$ & 国债对冲矢量 \\
$\mathcal{F}_i$ & 等于 $\mathcal{F}(t,t_F,T_i)$ \\
$X(t,T)$, $a(t_*,T)$ & $B(r;t;T)$ 表达式中出现的函数 \\
\bottomrule
\end{tabular}
\end{center}

\section*{第 10 章（远期利率的场论哈密顿量）}

\begin{center}
\begin{tabular}{@{}L{3.2cm}L{10.6cm}@{}}
\toprule
$f_{n,l}$ & 离散时间 $n$、未来时间 $l$ 的远期利率 \\
$F_n$ & 去掉最早远期利率 $f_{n,n}$ 的远期利率 \\
$F_{n+1}$ & 去掉最新远期利率的远期利率 \\
$S(n)$ & $t=n\epsilon$ 时刻远期利率作用量 \\
$\mathcal{L}_n[F_n;F_{n+1}]$ & 连接相邻两时间切片的远期利率拉格朗日量 \\
$\mathcal{V}_n$ & 远期利率在时刻 $n$ 的含时态空间 \\
$\mathcal{V}_{\mathrm{dual},n}$ & 时刻 $n$ 的含时对偶态空间 \\
$\mathcal{F}_n$ & 时刻 $n$ 远期利率物理态空间 \\
$\mathcal{H}_n$ & 时刻 $n$ 远期利率含时哈密顿量 \\
$\mathcal{I}_n$ & 离散时刻 $n$ 态空间的恒等算符 \\
$\vert f_t\rangle$, $\vert F_t\rangle$ & 时刻 $t$ 远期利率连续态空间基矢 \\
$\vert P(t,T)\rangle$, $\vert\mathcal{B}(t)\rangle$ & 零息/附息国债态矢量 \\
$\mathcal{H}_f(t)$, $\mathcal{H}_{f,z}(t)$ & 线性远期利率哈密顿量 \\
$\mathcal{H}_{\phi}(t)$ & 非线性远期利率哈密顿量 \\
$\mathcal{M}(x,x';t)$ & 随机波动率拉格朗日量的动能算符 \\
$\mathcal{H}(t)$ & 随机波动率线性远期利率哈密顿量 \\
$\mathcal{H}_{\mathrm{eff}}(t)$ & 有效远期利率线性厄米哈密顿量 \\
$\mathcal{O}(t)$, $\beta(t,x)$ & 确定 $\mathcal{H}_{\mathrm{eff}}(t)$ 所需的函数 \\
$\langle f_*\vert$ & 未来时刻 $t_*$ 远期利率对偶空间基 \\
$\mathcal{H}_f^*(t)\vert_{\mathcal{R}}$ & 限制在区域 $\mathcal{R}$ 上的远期利率哈密顿量 \\
$W$ & 时间积分的远期利率线性哈密顿量 \\
\bottomrule
\end{tabular}
\end{center}
